\documentclass[11pt,a4paper]{coverletter}

\senderblock
  {Radheem Bin Razi}
  {Distributed Systems \& Cloud-Native Engineer}
  {Ilmenau, Germany \textbar\ \href{mailto:sheikh.radheem@gmail.com}{sheikh.radheem@gmail.com} \textbar\ \href{https://radheem.github.io/my-cv}{portfolio}}

\begin{document}

%% ===== ENGLISH =====
\recipient{Haystack}{Hiring Team}{}

\opening{Dear Hiring Team,}

Building agents that can verify code and the mechanisms that judge model output is exactly the work I have been doing on my own time. In cv-tailor I deliberately split the system so that a pure, unit-tested ranker makes every decision while the LLM is confined to writing prose — the kind of verification boundary your role describes, where automated checks sit around model-generated code rather than trusting it. The combination of evaluating AI-generated code and designing verification for software-engineering tasks is the part of this posting I am most drawn to, and the part I am already practiced at.

For two years at Bluefin Exchange I was the design authority and system expert for the core services of a high-throughput crypto exchange. That meant reviewing other engineers' code and system designs daily, assessing code quality, and writing clear rationales for why a design or change was sound — then mentoring juniors through the same judgement. Curating correct code, spotting error patterns, and articulating *why* one solution is better than another is the muscle that role built, and it maps directly onto curating examples and evaluating model output for training.

The work spans the languages you list where I am strongest: Python, JavaScript with ReactJS, and Go. I have shipped full-stack, production-grade services across all three — Go microservices over gRPC and NATS in my IRS platform, Python with local llama.cpp and pgvector in a document-RAG system, and ReactJS/NodeJS backends at Bluefin and Al Hilal Invest. Deep familiarity with architecture, debugging and code review across these stacks is what lets me reason about whether AI-generated code is genuinely efficient, scalable and reliable.

I am currently a Master of Research student at TU Ilmenau and available now for the flexible, fully remote commitment you describe — comfortable anywhere from 10 to 40 hours a week with partial PST overlap. I would welcome the chance to help build the datasets and verification that make frontier coding models more reliable.

\closing{Sincerely,}{Radheem Bin Razi}

\clearpage
\selectlanguage{ngerman}

%% ===== DEUTSCH =====
\recipient{Haystack}{Personalabteilung}{}

\opening{Sehr geehrtes Hiring-Team,}

Agenten zu entwickeln, die Code verifizieren können, und die Mechanismen zu gestalten, die Modellausgaben bewerten, ist genau die Arbeit, mit der ich mich in meiner Freizeit beschäftige. Bei cv-tailor habe ich das System bewusst so aufgeteilt, dass ein reiner, unit-getesteter Ranker jede Entscheidung trifft, während das LLM auf das Verfassen von Prosa beschränkt bleibt — genau die Verifizierungsgrenze, die Ihre Stelle beschreibt, bei der automatisierte Prüfungen den modellgenerierten Code umgeben, anstatt ihm zu vertrauen. Die Kombination aus der Bewertung KI-generierten Codes und dem Entwurf von Verifizierung für Aufgaben der Softwareentwicklung ist der Teil dieser Ausschreibung, der mich am meisten anzieht — und der Teil, in dem ich bereits geübt bin.

Zwei Jahre lang war ich bei Bluefin Exchange der Design-Verantwortliche und Systemexperte für die Kerndienste einer durchsatzstarken Krypto-Börse. Das bedeutete, täglich den Code und das Systemdesign anderer Ingenieure zu überprüfen, die Codequalität zu bewerten und klare Begründungen dafür zu verfassen, warum ein Design oder eine Änderung tragfähig war — und anschließend Nachwuchskräfte durch dieselbe Urteilsbildung zu begleiten. Korrekten Code zu kuratieren, Fehlermuster zu erkennen und zu artikulieren, *warum* eine Lösung besser ist als eine andere, ist die Fähigkeit, die diese Rolle aufgebaut hat, und sie überträgt sich unmittelbar auf das Kuratieren von Beispielen und das Bewerten von Modellausgaben für das Training.

Die Arbeit erstreckt sich über die von Ihnen aufgeführten Sprachen, in denen ich am stärksten bin: Python, JavaScript mit ReactJS und Go. Ich habe in allen dreien produktionsreife Full-Stack-Dienste ausgeliefert — Go-Microservices über gRPC und NATS in meiner IRS-Plattform, Python mit lokalem llama.cpp und pgvector in einem Document-RAG-System sowie ReactJS/NodeJS-Backends bei Bluefin und Al Hilal Invest. Tiefe Vertrautheit mit Architektur, Debugging und Code-Review über diese Stacks hinweg ist es, was mir ermöglicht, zu beurteilen, ob KI-generierter Code wirklich effizient, skalierbar und zuverlässig ist.

Ich bin derzeit Master-of-Research-Student an der TU Ilmenau und ab sofort für das flexible, vollständig remote ausgeübte Engagement verfügbar, das Sie beschreiben — komfortabel zwischen 10 und 40 Stunden pro Woche mit teilweiser PST-Überlappung. Ich würde die Gelegenheit begrüßen, beim Aufbau der Datensätze und der Verifizierung mitzuwirken, die Coding-Modelle an der Spitze zuverlässiger machen.

\closing{Mit freundlichen Grüßen,}{Radheem Bin Razi}

\end{document}
